\clearpage
\section{계산 절차와 장 요약}
\label{sec:norm-trace-workflow-review}

\begin{learninggoal}
주어진 정보에 따라 행렬, 최소다항식, 체매장, 탑 공식, 유한체 공식을 선택하는 계산 절차를 정리한다. 노름과 대각합의 역할을 비교하고 다음 장의 근호확장과 가해군으로 연결한다.
\end{learninggoal}

\subsection{어떤 계산법을 선택할 것인가}

\begin{center}
\renewcommand{\arraystretch}{1.25}
\begin{tabular}{>{\raggedright\arraybackslash}p{0.25\textwidth}>{\raggedright\arraybackslash}p{0.30\textwidth}>{\raggedright\arraybackslash}p{0.34\textwidth}}
\toprule
주어진 정보 & 우선 사용할 도구 & 핵심 공식 \\
\midrule
기저와 곱셈표 &
곱셈행렬 &
$\Tr(a)=\operatorname{tr}(M_a)$,
$\Nm(a)=\det(M_a)$ \\

최소다항식 &
계수 읽기 &
$\Tr(\alpha)=-c_1$,
$\Nm(\alpha)=(-1)^n c_n$ \\

원소가 작은 부분체 생성 &
거듭제곱 공식 &
$\Tr_{L/K}(a)=[L:K(a)]\Tr_{K(a)/K}(a)$ \\

분리확장 또는 갈루아군 &
켤레들의 합과 곱 &
$\Tr(a)=\sum\sigma(a)$,
$\Nm(a)=\prod\sigma(a)$ \\

중간체의 탑 &
추이성 &
$\Tr_{L/K}=\Tr_{E/K}\circ\Tr_{L/E}$,
$\Nm_{L/K}=\Nm_{E/K}\circ\Nm_{L/E}$ \\

유한체 &
프로베니우스 &
$\Tr(a)=\sum a^{q^j}$,
$\Nm(a)=a^{(q^n-1)/(q-1)}$ \\
\bottomrule
\end{tabular}
\end{center}

\subsection{검산 순서}

노름과 대각합 계산에서는 다음 검산이 유용하다.
\begin{enumerate}
  \item \textbf{차수 확인}: 사용한 특성다항식의 차수가 $[L:K]$인지 확인한다.
  \item \textbf{부호 확인}: 일계수 $n$차다항식의 상수항 $c_n$에 대해
  \[
    \Nm(\alpha)=(-1)^n c_n
  \]
  이다.
  \item \textbf{기준체 원소 확인}: $c\in K$이면
  \[
    \Tr(c)=nc,\qquad\Nm(c)=c^n.
  \]
  \item \textbf{곱셈성 확인}: 복잡한 원소가 곱으로 분해되면 노름을 따로 계산한다.
  \item \textbf{탑 검산}: 중간체가 있으면 두 단계 계산과 직접 계산을 비교한다.
  \item \textbf{켤레 검산}: 분리확장에서는 모든 켤레의 합과 곱을 직접 확인한다.
  \item \textbf{판별식 검산}: 기저의 대각합 행렬식이 영이 아닌지 확인한다.
\end{enumerate}

\begin{pitfall}
최소다항식과 원소의 $L/K$-특성다항식을 항상 같은 것으로 착각하면 안 된다.
두 다항식이 같은 것은 $L=K(a)$일 때이다. $f_{a,K}$를 $a$의 최소다항식이라 하면 일반적으로
\[
  \chi_{a,L/K}(X)
  =f_{a,K}(X)^{[L:K(a)]}.
\]
\end{pitfall}

\begin{chapterreview}
\begin{itemize}
  \item $a\in L$은 곱셈사상 $m_a$를 통해 $K$-선형사상이 된다.
  \item
  \[
    \Tr_{L/K}(a)=\Tr(m_a),
    \qquad
    \Nm_{L/K}(a)=\det(m_a).
  \]
  \item 대각합은 $K$-선형이고 노름은 곱셈적이다.
  \item $L=K(\alpha)$이면 $\alpha$의 특성다항식과 최소다항식이 같다.
  \item 임의의 $a\in L$에 대해
  \[
    \chi_{a,L/K}
    =f_{a,K}^{[L:K(a)]}.
  \]
  \item 서로 다른 캐릭터와 서로 다른 체매장은 선형독립이다.
  \item 분리차수 $r$, 비분리차수 $q$에 대해
  \[
    \Tr(a)=q\sum_{i=1}^r\sigma_i(a),
    \qquad
    \Nm(a)=\left(\prod_{i=1}^r\sigma_i(a)\right)^q.
  \]
  \item 중간체 $K\subseteq E\subseteq L$에 대해 노름과 대각합은 추이적이다.
  \item 비분리 유한확장에서는 대각합이 영사상이다.
  \item 유한확장이 분리일 필요충분조건은 대각합 쌍이 비퇴화하는 것이다.
  \item 분리확장의 모든 기저는 유일한 대각합 쌍대기저를 갖는다.
  \item 유한체 확장에서 노름과 대각합은 프로베니우스의 곱과 합이다.
  \item 유한체의 노름과 대각합은 모두 전사이다.
  \item 서로소인 확장차수와 근의 지수 사이에서는 새로운 근이 생기지 않는다.
\end{itemize}
\end{chapterreview}

\subsection*{복습 카드}

\begin{itemize}
  \item \textbf{노름의 정의는?}
  $m_a$의 행렬식.
  \item \textbf{대각합의 정의는?}
  $m_a$의 대각합.
  \item \textbf{단순확장에서 최소다항식의 계수는?}
  $X^n+c_1X^{n-1}+\cdots+c_n$이면
  $\Tr(\alpha)=-c_1$, $\Nm(\alpha)=(-1)^nc_n$.
  \item \textbf{갈루아 확장에서의 공식은?}
  모든 자동동형상들의 합과 곱.
  \item \textbf{탑 공식은?}
  상대 대각합·노름을 먼저 계산한 뒤 아래 체로 다시 적용한다.
  \item \textbf{분리성의 대각합 판정은?}
  대각합 쌍이 비퇴화할 필요충분조건이 분리성이다.
  \item \textbf{유한체의 노름 지수는?}
  $(q^n-1)/(q-1)$.
  \item \textbf{유한체 대각합의 핵 크기는?}
  $q^{n-1}$.
\end{itemize}
